\chapter*{引言：什么是量子科学计算}
\addcontentsline{toc}{chapter}{引言：什么是量子科学计算}
\markboth{引言：什么是量子科学计算}{引言：什么是量子科学计算}

\emph{量子科学计算}（quantum scientific computing）研究如何用量子计算机求解科学与
工程中反复出现的大规模数值问题：线性方程组、特征值与矩阵函数、常微分与偏微分方程、
哈密顿量的时间演化，以及优化与机器学习。它处在两门学科的交汇处——一侧是以数值线性
代数与数值分析为核心的\emph{科学计算}，另一侧是以量子比特、量子线路与量子算法为语言
的\emph{量子计算}。本讲义的目标，是把后者系统地组织成一套能够攻击前者的工具箱，并在
每一步都诚实地交代：加速从何而来，又在何处要付出代价。

\section*{为什么要用量子计算机做科学计算}

许多最重要的科学计算问题，其规模随系统的自由度\emph{指数}增长。一个含 $n$ 个自旋的
量子多体系统，其希尔伯特空间维数为 $2^n$；一个 $d$ 维的偏微分方程，网格离散后的未知
量个数随 $d$ 指数增长（所谓``维数灾难''）。经典计算机必须把这指数多的数逐一存储、
更新，很快便超出任何机器的内存与时间预算。

量子计算机的根本不同在于：$n$ 个量子比特的态\emph{本身}就是一个 $2^n$ 维的复单位
向量，
\[
  \ket\psi=\sum_{x\in\{0,1\}^n} c_x\ket x,\qquad \sum_x\abs{c_x}^2=1 .
\]
于是仅用 $n$ 个物理比特，就``隐式地''承载了指数多的振幅 $c_x$；一个作用其上的量子
门，等价于对这 $2^n$ 维向量施加一个结构化的酉变换。这正是人们期待量子计算能为科学
计算带来加速的物理根源：态空间的维度与门作用的并行性都是指数级的。

这一想法有着清晰的思想脉络。Feynman 在 1982 年提出用一台受控的量子系统去\emph{模拟}
另一台量子系统，以绕开经典模拟的指数壁垒；Shor 的大数分解算法展示了量子计算在数论
问题上的超多项式优势；而 Harrow--Hassidim--Lloyd（2009）关于线性方程组的算法，则把
人们的注意力引向了一个更广阔的领域——用量子计算机做\emph{数值线性代数}，这正是本讲义
第二部分的主题。

\section*{一个核心的张力：指数维度与受限读出}

然而，指数维度本身并不免费地转化为加速。这里有一个贯穿全课程的核心张力。
\begin{itemize}[leftmargin=1.4em]
  \item \emph{写入（输入）。} 把一个一般的经典 $2^n$ 维向量 $b$ 加载为量子态
        $\ket b\propto\sum_x b_x\ket x$，一般需要 $\bigo(2^n)$ 的代价——除非 $b$ 具有
        可高效计算的结构（本讲义第 6 讲专门讨论何种态可被廉价制备）。
  \item \emph{读出（输出）。} 对 $\ket\psi$ 做一次测量只能得到一个随机的比特串 $x$，
        其概率为 $\abs{c_x}^2$；我们\emph{无法}一次性读出全部 $2^n$ 个振幅。若要把某个
        振幅、范数或期望值 $\bra\psi O\ket\psi$ 估计到精度 $\epsilon$，还需重复测量或
        借助振幅估计。
\end{itemize}
因此，量子科学计算真正能发挥优势的，是这样一类问题：输入能被高效地\emph{制备}、或由
神谕隐式地给出；计算过程可在指数维空间里高效完成；而我们最终想要的输出是一个
\emph{标量}（如某个期望值、范数或本征值）或一个可\emph{采样}的分布，而非整条解向量。
把握这条``输入--计算--输出''的边界，是理解每一个量子算法之加速何在、以及其``附加条款''
的关键。举例而言，量子线性系统算法并不直接交给你解向量 $x=A^{-1}b$ 的每个分量，而是
让你以远低于经典的代价制备出 $\ket x\propto A^{-1}\ket b$，再从中采样或读出某个
$\bra x M\ket x$——这既是它的威力所在，也是它的适用范围所限。

\section*{统一的算法语言：相位估计、块编码与 QSVT}

本课程的一条主线，是展示绝大多数量子数值算法都可以由少数几个\emph{原语}搭建而成：
\begin{itemize}[leftmargin=1.4em]
  \item \emph{相位估计}（quantum phase estimation）：读取酉算符本征相位的基本工具，
        是本征值问题、能量估计与振幅估计的共同引擎。
  \item \emph{振幅放大与振幅估计}：以平方级的加速搜索``好''的分支，或估计其权重。
  \item \emph{块编码}（block encoding）：把一个一般的（非酉的）矩阵 $A$ 嵌入一个更大的
        酉算符的左上角，从而让量子计算机得以``施加'' $A$。这是量子数值线性代数统一的
        输入模型，其代价由一个``次归一化因子'' $\alpha$ 刻画。
  \item \emph{酉算符线性组合}（LCU）与\emph{量子比特化}（qubitization）：把块编码做
        加法、乘法与多项式组合。
  \item \emph{量子奇异值变换}（QSVT）：对一个块编码矩阵的奇异值施加几乎任意的多项式
        变换，从而用一套框架统一了哈密顿量模拟、矩阵求逆、投影与谱滤波等看似迥异的
        任务。
\end{itemize}
学会用这套语言去``读''一个算法，往往比记住算法本身更有价值：一旦某个问题被化约为
``对某个块编码矩阵施加某个多项式''，它的量子线路与复杂度便几乎自动给出。本讲义反复
演练这种化约。

\section*{本讲义的组织}

全书分为五个部分，对应二十讲：
\begin{description}[leftmargin=0em,style=nextline]
  \item[第一部分 · 基础（第 1--5 讲）] 量子计算的基本公设、量子线路与通用门组；量子
        信道、距离与统计；Grover 搜索与振幅放大；量子傅里叶变换与相位估计；相位估计的
        应用。
  \item[第二部分 · 量子数值线性代数（第 6--10 讲）] 输入模型与块编码；LCU 与块编码的
        算术；量子比特化与量子信号处理（QSP）；QSVT 与矩阵函数；量子线性系统（QLSP）。
  \item[第三部分 · 微分方程与模拟（第 11--14 讲）] 基于乘积公式的哈密顿量模拟；基于
        块编码的模拟方法；线性常微分/偏微分方程算法；开放系统与非线性 PDE。
  \item[第四部分 · 张量网络（第 15--17 讲）] 张量网络态（矩阵乘积态 MPS）；DMRG 与
        基态；张量网络的时间演化（TEBD/TDVP）。
  \item[第五部分 · 优化与量子机器学习（第 18--20 讲）] 变分算法与 QAOA；量子线性规划
        与半定规划；量子机器学习。
\end{description}
前两部分奠定``把矩阵搬进量子计算机''的全部手段，后三部分则把这些手段分别投向微分
方程、强关联多体物理与优化/学习这三大应用领域。各讲彼此衔接：第三至第五部分反复
调用第二部分的块编码与 QSVT。

\section*{能期待怎样的加速——以及不能}

对量子加速保持清醒，是本课程的另一条主线。有几条需要牢记的``附加条款''：
\begin{itemize}[leftmargin=1.4em]
  \item \emph{输入/输出瓶颈。} 若态制备或读出本身就需要 $\bigo(2^n)$，则再快的中间
        过程也无济于事。真正的端到端加速通常要求输入由高效神谕给出、输出为标量或
        样本。
  \item \emph{对条件数、稀疏度与谱间隙的依赖。} 量子线性系统算法的复杂度依赖于矩阵的
        条件数 $\kappa$ 与稀疏度 $s$；块编码的代价由其因子 $\alpha$ 决定；基态制备依赖
        于谱间隙。宣称``指数加速''时，必须说明这些参数如何进入复杂度。
  \item \emph{容错与近期（NISQ）之别。} 本讲义中许多算法（相位估计、QSVT、QLSP）以
        \emph{容错}量子计算机为前提；而变分算法与 QAOA 则是为含噪的近期设备设计的，
        二者在假设、目标与可信的加速论断上截然不同。
\end{itemize}
把一个``量子加速''的论断拆解为输入、计算、输出三段，并逐段检查其隐含的代价，是我们
评判任何算法时反复演练的习惯。

\section*{预备知识与阅读建议}

阅读本讲义，假定读者熟悉线性代数（内积空间、谱定理、奇异值分解）、基本的数值分析
（误差、条件数、迭代法），以及初步的量子力学或量子信息（态、测量、酉演化）；\emph{不}
假定此前修读过量子计算课程——所需的一切将在第一部分从头建立。每一讲都以一个
\emph{学习目标}方框开篇，中间以定义、定理与\emph{带有完整线路和数值核验的例子}展开，
并以习题与延伸阅读收尾。建议初读时先抓住每讲的主线原语及其``输入--计算--输出''结构，
再回头推敲证明与例子的细节。读毕全书，你应当能够拿到一个具体的科学计算问题，判断它
是否适合量子方法、该用哪些原语搭建算法，并对所得加速给出有依据而非夸大的评估。
